\pagestyle{empty}

\cleardoublepage

% Capa do verso do livro
\begin{center}
\begin{tikzpicture}[remember picture,overlay]
    
    % Fundo preto
    \fill[bgColor] (current page.south west) rectangle (current page.north east);    

    % % Padrão de círculos central
    % \begin{scope}[shift={(current page.center)}, scale=1.0871]
    %     \foreach \x in {-20, -18, -16, -14, -12, -10, -8, -6, -4, -2, 0, 2, 4, 6, 8, 10, 12, 14, 16, 18} {
    %         \foreach \y in {-2, 2} {
    %             \begin{scope}[shift={(\x + 1, \y)}]
    %                 \foreach \r in {0.2, 0.4, 0.6, 0.8, 1.0} {
    %                     \draw[thick, lightgray] (0,0) circle (\r);
    %                 }
    %             \end{scope}
    %         }
    %         \foreach \y in {-3, -1, 1, 3} {
    %             \begin{scope}[shift={(\x, \y)}]
    %                 \foreach \r in {0.2, 0.4} {
    %                     \draw[thick, lightgray] (0,0) circle (\r);
    %                 }
    %             \end{scope}
    %         }
    %         \foreach \y in {-1, 1} {
    %             \begin{scope}[shift={(\x, 0)}]
    %                 \foreach \r in {0.2, 0.4, 0.6, 0.8, 1.0} {
    %                     \draw[thick, lightgray] (0,0) circle (\r);
    %                 }
    %             \end{scope}
    %         }
    %         \foreach \y in {-1, 1} {
    %             \begin{scope}[shift={(\x, \y)}]
    %                 \foreach \r in {0.2, 0.4, 0.6, 0.8, 1.0} {
    %                     \fill[bgColor] (0,0) circle (\r);
    %                 }
    %             \end{scope}
    %         }
    %         \foreach \y in {-1, 1} {
    %             \begin{scope}[shift={(\x, \y)}]
    %                 \foreach \r in {0.2, 0.4, 0.6, 0.8, 1.0} {
    %                     \draw[thick, lightgray] (0,0) circle (\r);
    %                 }
    %             \end{scope}
    %         }
    %     }
    % \end{scope}

    % Título
    \node[mainColor, anchor=north, text width =0.9\linewidth, align = justify] at ([yshift=-48pt] current page.north) {\scshape{Nunca é tarde demais para aprender ou revisar conceitos e ideias que nunca tivemos acesso antes, ou que estamos a muito tempo sem ter contato. Precisamos sempre ter a humildade de entender que subestimar o óbvio muitas das vezes é a causa principal da nossa frustração com conceitos que consideramos mais avançados.
    
    Esta apostila foi elaborada para tentar oferecer uma base mais sólida em matemática básica e dar um pouco mais de valor ao óbvio.}};

    % Logotipo da editora
    % \node[anchor=north, above=4pt of publisher.north] {\includegraphics[width=0.07\textwidth]{frontmatter/logo-white.png}};

\end{tikzpicture}
\end{center}